%
% halbkreiskontur.tex -- Halbkreiskontur fuer Hauptwerte
%
% (c) 2026 Damien Flury, OST Ostschweizer Fachhochschule
%
\documentclass[tikz]{standalone}
\usepackage{amsmath}
\usepackage{times}
\usepackage{txfonts}
\usetikzlibrary{arrows,math,decorations.markings}
\definecolor{darkblue}{rgb}{0,0,0.7}
\definecolor{darkred}{rgb}{0.7,0,0}
\begin{document}
\def\skala{1}
\begin{tikzpicture}[>=latex,thick,scale=\skala,
	midarrow/.style={
		postaction={decorate},
		decoration={markings,mark=at position #1 with {\arrow[line width=1.2pt]{latex}}}}]

% Parameter der Zeichnung
\def\R{3.6}
\def\x{1.4}
\def\e{0.75}

% Achsen
\draw[->] (-\R-0.6,0) -- (\R+0.7,0) coordinate[label={$\operatorname{Re}z$}];
\draw[->] (0,-0.6) -- (0,\R+0.7) coordinate[label={right:$\operatorname{Im}z$}];

% grosser Halbkreis C_R, im positiven Sinn durchlaufen
\draw[color=darkblue,line width=1.4pt,midarrow=0.5,midarrow=0.85]
	(\R,0) arc[start angle=0, end angle=180, radius=\R];
\node[color=darkblue] at (135:\R+0.42) {$C_R$};

% Strecke von -R bis x_0 - epsilon
\draw[color=darkblue,line width=1.4pt,midarrow=0.55]
	(-\R,0) -- (\x-\e,0);

% Einbuchtung C_epsilon: kleiner Halbkreis oberhalb der Polstelle,
% im negativen Sinn durchlaufen
\draw[color=darkblue,line width=1.4pt,midarrow=0.5]
	(\x-\e,0) arc[start angle=180, end angle=0, radius=\e];
\node[color=darkblue] at (\x,\e+0.34) {$C_\varepsilon$};

% Parametrisierung z(theta) = x_0 + epsilon e^{i theta}
\path (\x,0) ++(50:\e) coordinate (ztheta);
\draw[color=darkred,->] (\x,0) --
	node[color=darkred,font=\small,above left=-2pt] {$\varepsilon$} (ztheta);
\draw[color=darkred] (\x+0.32,0) arc[start angle=0, end angle=50, radius=0.32];
\path (\x,0) ++(25:0.5) coordinate (thetalabel);
\node[color=darkred,font=\small] at (thetalabel) {$\theta$};
\fill[color=darkred] (ztheta) circle[radius=0.05];
\node[color=darkred,font=\small] at (ztheta) [above right=-1pt] {$z(\theta)$};

% Strecke von x_0 + epsilon bis R
\draw[color=darkblue,line width=1.4pt,midarrow=0.55]
	(\x+\e,0) -- (\R,0);

% Polstelle x_0 auf der reellen Achse
\fill[color=darkred] (\x,0) circle[radius=0.06];
\node[color=darkred] at (\x,-0.42) {$x_0$};

% Pol z_1 in der oberen Halbebene
\fill[color=darkred] (-1.5,1.9) circle[radius=0.06];
\node[color=darkred] at (-1.5,1.9) [above right] {$z_1$};

% Beschriftung der Integrationsgrenzen
\draw (-\R,-0.08) -- (-\R,0.08);
\node at (-\R,-0.42) {$-R$};
\draw (\R,-0.08) -- (\R,0.08);
\node at (\R,-0.42) {$R$};
\draw[gray] (\x-\e,-0.08) -- (\x-\e,0.08);
\node[gray,font=\small] at (\x-\e,-0.92) {$x_0-\varepsilon$};
\draw[gray] (\x-\e,-0.18) -- (\x-\e,-0.72);
\draw[gray] (\x+\e,-0.08) -- (\x+\e,0.08);
\node[gray,font=\small] at (\x+\e,-0.92) {$x_0+\varepsilon$};
\draw[gray] (\x+\e,-0.18) -- (\x+\e,-0.72);

\end{tikzpicture}
\end{document}
